%
%

%%
\documentclass[12 pt, oneside]{book}
\usepackage{graphicx, fancyhdr, amsmath, times,  enumerate,geometry,makeidx,setspace,nomencl,eso-pic,xcolor,lipsum,calc,pst-node,tikz,fancybox,background,caption,subcaption,amsfonts}


\usetikzlibrary{calc}
\SetBgScale{1}
\SetBgAngle{0}
\SetBgColor{brown}
\SetBgOpacity{1}
%%
\geometry{verbose,a4paper,tmargin=1 in,bmargin=1in,lmargin=1.5 in,rmargin=.9 in}
%%%%%%%%%%%%%%%%%%%%%%%%%%%%%%%

\makeindex

%%
\usepackage[final]{pdfpages}
%%
%\setlength{\oddsidemargin}{2 cm}
%%
\newcommand{\VAtitle}[1]%
{\def\vtitle{#1}}%
\newcommand{\VAauthor}[1]%
{\def\vauthor{#1}}%
\newcommand{\VAadmissionyear}[1]%
{\def\vadmissionyear{#1}}%
\newcommand{\VAacademicyear}[1]%
{\def\vacademicyear{#1}}%
\newcommand{\VAregisternumber}[1]%
{\def\vregisternumber{#1}}%
\newcommand{\VAprincipal}[1]%
{\def\vprincipal{#1}}%
\newcommand{\VAguide}[1]%
{\def\vguide{#1}}%
\newcommand{\VAguidedg}[1]%
{\def\vguidedg{#1}}%
\newcommand{\VAhod}[1]%
{\def\vhod{#1}}%
\newcommand{\VAmonth}[1]%
{\def\vmonth{#1}}%
\newcommand{\VAdept}[1]%
{\def\vdept{#1}}%
\newcommand{\VAclass}[1]%
{\def\vclass{#1}}%
\newcommand{\VApaper}[1]%
{\def\vpaper{#1}}%
%%
%%
%\renewcommand\bibname{References}




\SetBgContents{}

\VApaper{SEMINAR REPORT}


\usepackage{titlesec}
\titleformat{\chapter}[display]
{\normalfont\huge\bfseries\centering}{\chaptertitlename\ \thechapter}{20pt}{\Huge}

\begin{document}
	
\VAtitle{Seminar Report Title}
%
%   In the next line, replace "Student" by your name. 
%   Write full name without, Mr. or Ms.
%
\VAauthor{Student Name}%
%
%
\VAadmissionyear{2017}%
%
\VAacademicyear{2020-2021}%


\VAregisternumber{University Register Number}% 


\VAprincipal{Principal Name}% 
%	Name of your Mentor
\VAguide{Guide Name}% 
%	Designation and Department of the Mentor
\VAguidedg{Asst. Prof.,} 
\VAhod{HoD Name}
%
%   Enter the Sem End month like "December"
%   Type Month and Year
%
\VAmonth{November 2018}%
\VAdept{Branch of Engineering}%
\VAclass{Seventh Semester B.Tech}%

%
%  
\fancypagestyle{plain}{%
\fancyhf{} % clear all header and footer fields
\fancyhead[L]{{\scriptsize \vtitle}}
\fancyhead[R]{\includegraphics[width=0.5cm]{layout_files/VidyaLogo.JPG}}
%\fancyfoot[C]{\bfseries \thepage} 
\fancyfoot[L]{{\footnotesize Department of Electrical \& Electronics Engg.}}
\fancyfoot[R]{\footnotesize VAST, Thalakotukara}
\fancyfoot[C]{\footnotesize \bf \thepage}%
\renewcommand{\headrulewidth}{1pt}%
\renewcommand{\footrulewidth}{1pt}%
}%
%
%
\pagestyle{empty}
%

\begin{spacing}{1.5}
\begin{titlepage}


\SetBgContents{
\begin{tikzpicture}[overlay,remember picture]
\draw [line width=3pt]
    ($ (current page.north west) + (3.0cm,-1.8cm) $)
    rectangle
    ($ (current page.south east) + (-1.35cm,1.8cm) $);
\draw [line width=1pt]
    ($ (current page.north west) + (3.15cm,-1.95cm) $)
    rectangle
    ($ (current page.south east) + (-1.5cm,1.95cm) $); 
\end{tikzpicture}
}

\begin{center}


{ \LARGE \rmfamily \bf \vtitle}\\[1 cm]

{ \large \rmfamily \vpaper \\ SUBMITTED IN PARTIAL FULFILLMENT OF THE \\REQUIREMENTS FOR THE AWARD OF DEGREE OF}\\[1 cm]

{ \Large \rmfamily \bf BACHELOR OF TECHNOLOGY}\\
{\large \rmfamily in}\\[.4 cm]
{ \Large \rmfamily \bf ELECTRICAL AND ELECTRONICS ENGINEERING}\\
{\large \rmfamily of}\\[.4 cm]
{ \Large \rmfamily \bf APJ Abdul Kalam Technological University}\\
{\large \rmfamily by}\\[1 cm]
{\large \rmfamily \bf \vauthor} \\
{\large \bf Univ. Reg.No. \vregisternumber} \\[2 cm]
%

%
\includegraphics[width=3.5 cm]%
{layout_files/VidyaLogo.JPG}\\
\scriptsize (AN ISO 9001:2008 CERTIFIED INSTITUTION )\\[1.5 cm]

%
{\Large \bf Department of \vdept} \\
{\Large \rmfamily Vidya Academy of Science \& Technology\\[.2 cm]
\large Thalakkottukara, Thrissur - 680 501}\\
({ \bf \tt http://www.vidyaacademy.ac.in})\\[.1 cm]
{\large \rmfamily \vmonth}


\end{center}
\end{titlepage}
%
%******************************************************
%
%
\clearpage
\pagenumbering{roman}
%


%\addcontentsline{toc}{chapter}{\quad Certificate}

\begin{titlepage}
	
\SetBgContents{
\begin{tikzpicture}[remember picture,overlay]

\draw [line width=3pt]
    ($ (current page.north west) + (3 cm,-1.8cm) $)
    rectangle
    ($ (current page.south east) + (-1.5cm,1.8cm) $);
\draw [line width=1pt]
    ($ (current page.north west) + (3.15cm,-1.95cm) $)
	rectangle
    ($ (current page.south east) + (-1.65cm,1.95cm) $); 
\end{tikzpicture}
}


\begin{center}


{\Large \bf Department of \vdept  }\\
{\Large \bf Vidya Academy of Science \& Technology}\\
{\normalsize \bf Thalakkottukara, Thrissur - 680 501\\
({\tt http://www.vidyaacademy.ac.in})}\\[0.75cm]
%
%   Logo
%
\includegraphics[width=3.5 cm]{layout_files/VidyaLogo.JPG}\\
\scriptsize (AN ISO 9001:2008 CERTIFIED INSTITUTION )\\[1 cm]
%
 \Huge  $ \mathfrak{Certificate}$\\[0.5cm]
%
\end{center}
\index{Certificate}
\index{University of Calicut}

\quad This is to certify that the seminar report titled {\bf ``\vtitle"} is a bonafide record of the work carried out by {\bf \vauthor\ } {\bf (Univ. Reg.No. \vregisternumber) } of Vidya Academy of Science \& Technology, Thalakkottukara, Thrissur - 680 501 in partial fulfillment of the requirements for the award of {\bf Degree of Bachelor of Technology} in {\bf Electrical and Electronics Engineering} of  {\bf APJ Abdul Kalam Technological University}, during the academic year \vacademicyear. \\[2 cm]
 
\noindent{\bf Seminar Guide/Supervisor} \hfill  {\bf Head of Department} \\[.3cm]
\noindent \vguide \hfill \vhod \\ \vguidedg\ Dept. of EEE \hfill Prof., Dept. of ***  



%
\end{titlepage}

%   End of Certificate
%  
%
\clearpage
\pagestyle{plain}
%

\addcontentsline{toc}{chapter}{\quad ACKNOWLEDGEMENT}

\chapter*{Acknowledgement}
%


\par
\hspace{0.9cm}I wish to record my indebtedness and thankfulness 
to all those who helped me prepare this report titled ``{\bf \vtitle\ }''  and present it in a satisfactory way.

\vspace {.2cm}
\par
\hspace{.35cm}First and foremost I thank God Almighty for His providence and for being the guiding light throughout the seminar.

\vspace{0.2cm}
\par
\hspace{0.35cm}I would like to thank my guide {\bf \vguide, } \vguidedg\ of \vdept\  Dept. for providing critical inputs in the preparation of this report. I also thank all other faculty members in our department for their guidance.


\vspace{.2cm}
\par 
\hspace{.35cm}I am thankful to {\bf \vhod}, 
Head of \vdept\  Department, and our Principal {\bf \vprincipal}, for their sole co-operation.



\vspace{0.2cm}
\par
\hspace{0.35cm}Finally, I would like to extend my sincere gratitude to friends who have always been helpful, in preparing and presenting the report and in the discussion following the presentation.

\begin{flushright}
\vauthor\\
Reg. No. \vregisternumber\\
\vclass\  ( \vadmissionyear\  Admissions)\\
Vidya Academy of Science \& Technology\\
Thrissur - 680 501.
\end{flushright}

\vmonth\ 


%
\clearpage
%
	% Dont Edit . . 



\chapter*{Abstract}
\addcontentsline{toc}{chapter}{\quad ABSTRACT}

Copy your Abstract of the seminar here . . .

%%%%%%%%%%%%%%%%%%%%%%%%%%%%%%%%%%%%
%%%%%%%%%%%%%%%%%%%%%%%%%%%%%%%%%%%%
\tableofcontents
\addcontentsline{toc}{chapter}{\quad  LIST OF FIGURES}
\listoffigures
\addcontentsline{toc}{chapter}{\quad  LIST OF  TABLES}
\listoftables
% For adding List of symbols or abbreviations
\addcontentsline{toc}{chapter}{\quad  LIST OF SYMBOLS AND ABBREVIATIONS}

\chapter*{List of Symbols and Abbreviations}


\begin{tabbing}


\hspace{1cm}\= {$ \bf VAST  $}\quad\= Vidya Academy of Science and Technology\\[5pt]

\> {$ \bf A_i  $} \> Area of the $i^{th}$ component\\[5pt]

\> {$ \bf DC $} \>   Direct current electricity\\[5pt]

\> {$ \bf DSO  $} \> Digital Storage Oscilloscope\\[5pt]

\> {$ \bf FBI  $} \> Full Bridge Inverter


\end{tabbing}
\mainmatter			
%%%%%%%%%%%%%%%%%%%%%%%%%%%%%%%%%%%%
%%%%%%%%%%%%%%%%%%%%%%%%%%%%%%%%%%%%
%   The main contents of the paper begin here. 
%   In the following replace with 
%   the title of the first chapter of your paper.
%
\chapter{INTRODUCTION}

{\em Give some intro abt this chapter . . .. .  . }

\section{General}    % For giving Section  eg: 1.1

Give some introduction abt the project, maximum 1 or 2 pages   . . .  

Electrical engineering is a field of engineering that generally deals with the study and application of electricity, electronics, and electromagnetism. This field first became an identifiable occupation in the latter half of the 19th century after commercialization of the electric telegraph, the telephone, and Fig. \ref{a}
%%%%%%%%%%%%%%%%%%%%%%%%%%%%%%%%%%%
%		Before using image put the jpg image file in the folder. 
%		Here   1.jpg   is the file name
%%%%%%%%%%%%%%%%%%%%%%%%%%%%%%%%%%%
\begin{figure}[h]
	\begin{center}
		\includegraphics[width =5 cm]{1.jpg}
		\caption{Resonance based inertial Electromagnetic Microgenerator}
		\label{ab}
	\end{center}
\end{figure}

\begin{figure}[h]
\begin{center}
\includegraphics[width =5 cm]{1.jpg}
\caption{Resonance based inertial Electromagnetic Microgenerator}
\label{a}
\end{center}
\end{figure}

electric power distribution and use.\cite{ref1} Photovoltaic (PV) \index{Photovoltaic}power supplied to the system is gaining more and more visibility. \cite{ref3}

Photovoltaic (PV) power supplied to the system is gaining more and more visibility Solid-state inverters have shown for putting PV systems now. \cite{ref2}



 \section{Objectives of the Work}     % You can add any sections like that . . 

The concept of micro-inverter (also known as a module integrated converter / inverter) has . . . 

A plug-in DMRC system under consideration, where R(z) is the reference input, Y (z) is the output, E(z) = R(z) − Y (z) is the tracking error, D(z) is the disturbance, Gc(z) is the conventional feedback controller, Gs(z) is the plant, Gdr(z) is . . . . 

\section{Outline of Report}

First Chapter contains  . . . . .. .. .. Second chapter contains . . .



\chapter{TOPOLOGIES}			%Difft Sections in IEEE Paper

{\em Give some intro abt this chapter . . .. .  .

A plug-in DMRC system under consideration, where R(z) is the reference input, Y (z) is the output, . . .  }

\section{PRIMARY SIDE SWITCH TOPOLOGY}

A plug-in DMRC system under consideration, where R(z) is the reference input, Y (z) is the output, E(z) = R(z) − Y (z) is the tracking error, D(z) is the disturbance,
Gc(z) is the conventional feedback controller, Gs(z) is the plant, Gdr(z) is a feedforward plug-in \cite{ref3}.


\section{SECONDARY SIDE SWITCH TOPOLOGY}

A plug-in DMRC system under consideration, where R(z) is the reference input, Y (z) is the output, E(z) = R(z) − Y (z) is the tracking error, D(z) is the disturbance,

\begin{figure}[h]
\begin{center}
\includegraphics[width = 8 cm]{5.jpg}
\caption{Operation modes of the split capacitor topology converter}
\end{center}
\end{figure}

A plug-in DMRC system under consideration,
where R(z) is the reference input, Y (z) is the output, E(z) =
R(z) − Y (z) 

\begin{equation}
d_fT_s=i_{pk}/m_2=i_{pk}L_1/(V_0-v_{ik})
\end{equation}

During this switching cycle, the energy ($E_k$) transferred
from the input to the output can be obtained as:

\begin{equation}
E_k=v_{ik}.i_{pk}.T_s(D+d_f)/2
\end{equation}


\section{SECONDARY SIDE SWITCH TOPOLOGY}

A plug-in DMRC system under consideration,
where R(z) is the reference input, Y (z) is the output, . . .

For equ..

\begin{equation}
d_fT_s=i_{pk}/m_2=i_{pk}L_1/(V_0-v_{ik})
\end{equation}

During this switching cycle, the energy ($E_k$) transferred
from the input to the output can be obtained as:

\begin{equation}
E_k=v_{ik}.i_{pk}.T_s(D+d_f)/2
\end{equation}


\begin{figure}[h]
        \centering
        \begin{subfigure}[b]{0.3\textwidth}
                \includegraphics[width=\textwidth]{5}
                \caption{Circuit A}
                \label{fig:Circuit A}
        \end{subfigure}%
        ~ %add desired spacing between images, e. g. ~, \quad, \qquad, \hfill etc.
          %(or a blank line to force the subfigure onto a new line)
        \begin{subfigure}[b]{0.3\textwidth}
                \includegraphics[width=\textwidth]{5}
                \caption{Circuit B}
                \label{fig:Circuit B}
        \end{subfigure}
        ~ %add desired spacing between images, e. g. ~, \quad, \qquad, \hfill etc.
          %(or a blank line to force the subfigure onto a new line)
        \begin{subfigure}[b]{0.3\textwidth}
                \includegraphics[width=\textwidth]{5}
                \caption{Circuit C}
                \label{fig:Circuit C}
        \end{subfigure}
        \caption{Circuit A, B and C}\label{fig:Circuits}
\end{figure}



\chapter{ RESULTS OF THE SETUP}


A plug-in DMRC system under consideration, where R(z) is the reference input, Y (z) is the output, E(z) = R(z) − Y (z) is the tracking error, D(z) is the disturbance, \.{P}
Gc(z) is the conventional . . .


To use bullets

\begin{itemize}
\item{Presented}
\item{Paper  }
\item{Presented}
\item{Paper  }
\end{itemize}

General, a micro-inverter PV \index{Photovoltaic} is often provided by a low voltage solar panel, which requires a high voltage lift ratio to produce the desired AC output .


\begin{itemize}
\item{Presented}
\item{Paper  }
	\begin{itemize}
		\item{Presented}
		\item{Paper  }
	\end{itemize}
\item{Presented}
\item{Paper  }
\end{itemize}


Inverter has become a future trend for single phase photovoltaic (PV) power. In general, a micro-inverter PV \index{Photovoltaic} is often provided by a low voltage solar panel . . . \vspace{2 cm} % Vspace to leave 2 cm space

To use 1,2,3, numbers

\begin{enumerate}
\item prepare a source file with the extension "tex"
\item Compile it with \LaTeX to produce a "dvi" file
\item Print the document using a "dvi" driver
\end{enumerate}


For making a table . .. .

\begin{table}[h]
\centering
\begin{tabular}{|l|c|}
\hline
Planet & Diameter(km) \\
\hline
Mercury & 4878\\
Pluto & 2274\\
Mercury & 4878\\
Pluto & 2274\\
Mercury & 4878\\
Pluto & 2274\\
Mercury & 4878\\
Pluto & 2274\\
\hline
\end{tabular}
\caption{Table Name}
\end{table}




\chapter{CONCLUSION}


Controller, Gf (z) is a filter to obtain a stable overall closed-loop system, and Q(z) is a lowpass filter to enhance the robustness of the overall system at the cost of imperfect elimination of high-frequency errors.
The feedback controller Gc(z) is chosen so that the transfer function.



%
%
%
% ***   Creating Bibliography   ***
%
% Do not delete the following two lines
%
\clearpage
\addcontentsline{toc}{chapter}{\quad BIBLIOGRAPHY}
%
%   The markups for creating the bibliography 
%   begins here. Do not change the two lines 
%   "\begin{thebibliography}{99}" and 
%   "\end{thebibliography}".
%   One sample bibliography item is included. 
%   Each new item is to be preceded by
%   "\bibitem{}".
%   Add additional items if there are any.
%   Bibliography styles:
%   Titles of books   : Italics  {use {\em Title} )
%   URLs of websites  : Type-writer (use {\tt Title} )
%   Titles of papers  : In double quotes (use ``Title")
%
\begin{thebibliography}{99}
%
\bibitem{ref1}
Suman Dwari, Leila Parsa, and Rohan Dayal, “Design and Implementation of a Direct AC–DC Boost Converter for Low-Voltage Energy Harvesting,” {\em IEEE Transactions on Industrial Electronics}, Vol. 58, No. 6, June 2011, pp. 2387–2396.

\bibitem{ref2}
Suman Dwari and Leila Parsa, "Efficient Low Voltage Direct AC/DC Converters for Selfpowered Wireless Sensor Nodes and Mobile Electronics”, {\em Telecommunications Energy Conference}, Intelec Sept. 2008.

\bibitem{ref3}
Suman Dwari and Leila Parsa,, "An Efficient AC–DC Step-Up Converter for Low-Voltage Energy Harvesting," {\em IEEE Transactions on Power Electronics}, Vol. 25, No. 8, August 2010,pp. 2188-2199.




%
\end{thebibliography}
%
%   ***   Creating Index   ***
%
%   To add a particular word, say abcxyz, to index write 
%   the command \index{abcxyz} near the word where it 
%   first appears in the text. Include all technical 
%   sounding words in the index. 
%   Do not delete the following three lines.
%
%\clearpage
%\addcontentsline{toc}{chapter}{\quad INDEX}
%\printindex
%
\clearpage
\addcontentsline{toc}{chapter}{\quad APPENDIX}
\chapter*{\centering{APPENDIX}}
%%%%%%%%%%%%%%%%%%%%%%%%%%%%%%%%%%%%%
%%                          
%%		APPENDIX
%%    if u have pgms make that a pdf & add it here, like the data sheets
%%
%%    If you need to give an Intro to the Appendix
%%                                 Otherwise Delete it . . . 
%%
%%%%%%%%%%%%%%%%%%%%%%%%%%%%%%%%%%%%%
%%%%%%%%%%%%%%%%%%%%%%%%%%%%%%%%%%%%%

\begin{itemize}
\item Reference Paper of the Seminar
\item PPT of the Seminar
\end{itemize}

%%%%%%%%%%%%%%%%%%%%%%%%%%%%%%%%%%%%%
%%%%%%%%%%%%%%%%%%%%%%%%%%%%%%%%%%%%%
%%                      
%%                      For Adding PDF (Data Sheets) for the Appendix
%%
%%%%%%%%%%%%%%%%%%%%%%%%%%%%%%%%%%%%%
%%%%%%%%%%%%%%%%%%%%%%%%%%%%%%%%%%%%%


\includepdf[pages=-]{paper.pdf}    % for adding Paper in PDF (same folder)
\includepdf[pages=-]{ppt.pdf}    % for adding PPT in PDF in 1 Page 6 Slides


%   Printing the last page
%
\input{layout_files/lastpage}
%
%
%   ***   The end   ***
%
%
\end{document}

